\chapter{جمل الفرائض: الزكاة والحج وأحكام النساء والطعام}

\section{جمل الفرائض في الزكاة وتفصيل السنة لها}
الحمد لله والصلاة والسلام على رسول الله، أما بعد؛ فقد شرع الإمام الشافعي رحمه الله تعالى في بيان جمل الفرائض التي نزل أصلها في القرآن الكريم وأتت تفصيلاتها وتفريعاتها في السنة النبوية المطهرة. وابتدأ بالزكاة بعد الصلاة.

الزكاة في اللغة: الطهارة والنماء. وفي الشرع: أخذ مال مخصوص من أشخاص مخصوصين في أوقات معلومة لأصناف معلومة.
قال الله جل وعلا: ﴿وَأَقِيمُوا الصَّلَاةَ وَآتُوا الزَّكَاةَ﴾، وقال في ذم الكفار: ﴿الَّذِينَ هُمْ يُرَاءُونَ * وَيَمْنَعُونَ الْمَاعُونَ﴾، فسر بعض السلف (كعلي وابن عباس رضي الله عنهم) الماعون هنا بالزكاة المفروضة. 

وقال تعالى: ﴿خُذْ مِنْ أَمْوَالِهِمْ صَدَقَةً تُطَهِّرُهُمْ وَتُزَكِّيهِمْ بِهَا﴾. مخرج هذه الآية عام، فيحتمل أن تؤخذ الزكاة من جميع الأموال، فجاءت السنة المطهرة فدلت على أن الزكاة في بعض الأموال دون بعض، وحددت الأنصبة والمقادير:

\subsection{زكاة الماشية}
لما كانت الماشية أصنافاً، أخذ رسول الله صلى الله عليه وسلم الزكاة من الإبل والبقر والغنم، ولم يأخذ من الخيل والبغال والحمير، وقال: «ليس على المسلم في عبده ولا فرسه صدقة». وحدد أنصابها؛ فجعل زكاة الإبل تبتدئ من خمس، والبقر من ثلاثين، والغنم من أربعين، فاستدللنا بسنته على وجوب الزكاة فيما أخذ منه وسقوطها عما ترك.

\subsection{زكاة الزروع والثمار}
أخذ رسول الله صلى الله عليه وسلم الزكاة من النخل والعنب بالخَرْص (وهو تقدير الخبير لما على الشجر من ثمر)، وجعل فيها العشر إن سقيت بماء السماء أو العيون، ونصف العشر إن سقيت بالدلو والمشقة. ولم يأخذ من باقي الشجر (كالتين والجوز واللوز).
كما أخذ الزكاة من الحنطة (القمح) والشعير والذرة وكل ما يقتات ويدخر، ولم يأخذ مما لا يقتات ويدخر (كحب الرشاد والكزبرة).

\subsection{زكاة النقدين والركاز}
وفرض صلى الله عليه وسلم الزكاة في الورِق (الفضة) والذهب لأنهما نقد الناس وثمن مبيعاتهم، ولم يأخذ من النحاس والحديد والرصاص شيئاً. وبين أن زكاة النقدين والماشية تؤخذ مرة كل حَول (سنة)، بينما زكاة الزروع تؤخذ يوم الحصاد لقوله تعالى: ﴿وَآتُوا حَقَّهُ يَوْمَ حَصَادِهِ﴾. وسنّ في الركاز (ما استخرج من دفائن الأرض) الخُمس في يوم خروجه.

\section{جمل الفرائض في الحج ومناسك السنة}
قال تعالى: ﴿وَلِلَّهِ عَلَى النَّاسِ حِجُّ الْبَيْتِ مَنِ اسْتَطَاعَ إِلَيْهِ سَبِيلًا﴾. فذكر الفرض مجملاً.
فبين رسول الله صلى الله عليه وسلم أن "السبيل" هو الزاد والراحلة، وعلمنا كيف نلبي، وما يلبس المحرم وما يجتنبه (كالطيب والثياب المخيطة)، وعلمنا أعمال عرفة ومزدلفة والرمي والحلق والطواف؛ مردداً: «خذوا عني مناسككم».
فلو أن امرأً لم يعلم من السنة إلا ما فصله النبي صلى الله عليه وسلم لمجمل القرآن لكانت الحجة قائمة عليه بوجوب طاعة رسول الله صلى الله عليه وسلم، ووجب عليه أن يعلم أن سنة رسول الله لا تخالف كتاب الله أبداً.

\section{أحكام العدد والمحرمات من النساء}

\subsection{عدة المتوفى عنها زوجها الحامل}
قال تعالى: ﴿وَالَّذِينَ يُتَوَفَّوْنَ مِنْكُمْ وَيَذَرُونَ أَزْوَاجًا يَتَرَبَّصْنَ بِأَنْفُسِهِنَّ أَرْبَعَةَ أَشْهُرٍ وَعَشْرًا﴾، وقال تعالى: ﴿وَأُولَاتُ الْأَحْمَالِ أَجَلُهُنَّ أَنْ يَضَعْنَ حَمْلَهُنَّ﴾.
فظن بعض أهل العلم أن الحامل المتوفى عنها زوجها تعتد بأبعد الأجلين (تجمع بين العدتين). فجاءت السنة لتبين الأمر؛ فلما وضعت سُبيعة الأسلمية رضي الله عنها حملها بعد وفاة زوجها بأيام قليلة، قال لها رسول الله صلى الله عليه وسلم: «قد حللت فتزوجي». فدلت السنة على أن وضع الحمل ينهي العدة مطلقاً وإن كان بعد سويعات من الوفاة.

\subsection{المحرمات من النساء (الجمع بين المرأة وعمتها)}
قال تعالى معدداً المحرمات في النكاح: ﴿حُرِّمَتْ عَلَيْكُمْ أُمَّهَاتُكُمْ وَبَنَاتُكُمْ...﴾ ثم قال: ﴿وَأُحِلَّ لَكُمْ مَا وَرَاءَ ذَلِكُمْ﴾.
فكان ظاهر الآية جواز نكاح أي امرأة لم تذكر في الآية. لكن السنة المطهرة بينت أن قوله ﴿وَأُحِلَّ لَكُمْ مَا وَرَاءَ ذَلِكُمْ﴾ مقيد بما لم تحرمه السنة أيضاً. فقد نهى النبي صلى الله عليه وسلم أن يجمع الرجل بين المرأة وعمتها، وبين المرأة وخالتها. والتحريم هنا ليس تحريماً مؤبداً لذات العمة أو الخالة كتحريم الأم والأخت، بل هو تحريم "الجمع" بينهما في عصمة واحدة فقط.

\subsection{حتى تنكح زوجاً غيره (اشتراط الوطء)}
قال تعالى في المطلقة ثلاثاً: ﴿فَإِنْ طَلَّقَهَا فَلَا تَحِلُّ لَهُ مِنْ بَعْدُ حَتَّى تَنْكِحَ زَوْجًا غَيْرَهُ﴾. احتمل لفظ "النكاح" العقد المجرد، واحتمل العقد مع الإصابة (الجماع). فجاءت السنة موضحة لمراد الله؛ حينما جاءت امرأة رفاعة القرظي تشتكي أن زوجها الثاني لم يصبها، فقال لها النبي صلى الله عليه وسلم: «أتريدين أن ترجعي إلى رفاعة؟ لا، حتى تذوقي عُسَيْلَتَهُ ويذوق عُسَيْلَتَكِ». فبينت السنة أن التحليل لا يتم إلا بنكاح صحيح يتبعه وطء.

\section{ما أحل وحرم من الطعام}
قال تعالى: ﴿قُلْ لَا أَجِدُ فِي مَا أُوحِيَ إِلَيَّ مُحَرَّمًا عَلَى طَاعِمٍ يَطْعَمُهُ إِلَّا أَنْ يَكُونَ مَيْتَةً أَوْ دَمًا مَسْفُوحًا أَوْ لَحْمَ خِنْزِيرٍ...﴾.
ظاهر الآية حصر المحرمات في هذه المذكورات، ولكن السنة المطهرة جاءت بتحريم "كل ذي ناب من السباع، وكل ذي مخلب من الطير". فدلت السنة إما على أن الحصر في الآية كان لما سألوا عنه حينها، أو أنه مخصوص، فما حرمه رسول الله فهو حرام كأنما نزل في القرآن.

\section{ما تجتنبه المعتدة للوفاة}
قال تعالى في المعتدة للوفاة: ﴿فَإِذَا بَلَغْنَ أَجَلَهُنَّ فَلَا جُنَاحَ عَلَيْكُمْ فِيمَا فَعَلْنَ فِي أَنْفُسِهِنَّ بِالْمَعْرُوفِ﴾.
القرآن فرض عليها أن تتربص أربعة أشهر وعشراً، ولم يذكر الأشياء التي تجتنبها المعتدة، فجاءت السنة النبوية فألزمتها بأمرين زائدين على مجرد التربص عن الأزواج: 
1. الامتناع عن الزينة والطيب وما يدعو للرغبة فيها.
2. لزوم السكن والمكث في بيت زوجها الذي مات وهي فيه، فلا تخرج منه إلا لحاجة ضرورية.

\section{خاتمة: العذر بالجهل لمن خالف السنة من العلماء}
يقرر الشافعي أصلاً منهجياً عظيماً: كل عالم رُوي عنه قول يخالف سنة صحيحة لرسول الله صلى الله عليه وسلم، فإنه معذور مأجور؛ لأنه إنما خالفها لجهله بها وعدم بلوغها إياه، ولو علمها لرجع إليها وترك قوله. ولا يجوز لمسلم إذا بلغته السنة أن يتركها لقول عالم مهما علا شأنه؛ فالحجة في قول رسول الله صلى الله عليه وسلم، ومن أطاع الرسول فقد أطاع الله.